% !TEX root = ../main.tex
%% Corrigé du devoir surveillé n° 4 — valeurs vérifiées : verifications/devoirs.py

\corrigeDevoir{4}

\begin{corrigeBac}{7}{La Station spatiale internationale}
\begin{questions}
\item $F_\mathrm{T/ISS}=\dfrac{G\,M_\mathrm{T}\,m}{(R_\mathrm{T}+h)^2}$, avec $R_\mathrm{T}+h=\qty{6,771e6}{m}$ :
$F=\dfrac{\num{6,67e-11}\times\num{5,97e24}\times\num{4,20e5}}{(\num{6,771e6})^2}=\res{\qty{3,65e6}{N}}$. \pts{1,5 pt}
\item Force attractive, dirigée de l'ISS vers le centre de la Terre, représentée par une flèche de \res{\qty{3,65}{cm}} (schéma). \pts{1 pt}
\item $g=\dfrac{F}{m}=\dfrac{G\,M_\mathrm{T}}{(R_\mathrm{T}+h)^2}=\res{\qty{8,69}{N.kg^{-1}}}$ : à peine \res{\qty{11}{\percent}} de moins qu'au sol. \pts{1,5 pt}
\item $F_\mathrm{L}=\dfrac{\num{6,67e-11}\times\num{7,35e22}\times\num{4,20e5}}{(\num{3,84e8})^2}=\res{\qty{14,0}{N}}$, environ $\res{\num{2,6e5}}$ fois plus faible que la
force exercée par la Terre : elle est négligeable. \pts{1,5 pt}
\item L'affirmation est \res{\text{fausse}} : à \qty{400}{km}, la Terre attire encore les astronautes presque autant qu'au sol. Ils flottent
parce que la station et eux sont en \res{\text{chute libre}} autour de la Terre : ils « tombent » ensemble, sans jamais atteindre
le sol, grâce à leur grande vitesse horizontale. \pts{1,5 pt}
\end{questions}
\schema[corrige]{ds4-iss}
\end{corrigeBac}

\begin{corrigeBac}{7}{La plongée sous-marine}
\begin{questions}
\item Les molécules d'eau, en mouvement désordonné, \res{\text{frappent}} sans cesse la paroi : la force pressante résulte de
l'ensemble de ces chocs. \pts{1 pt}
\item Les points sont alignés sur une droite : $P=P_0+k\,h$ avec $k\approx\qty{0,101}{bar.m^{-1}}=\qty{1,01e4}{Pa.m^{-1}}$ (graphique). C'est
la loi $P(h)=P_0+\rho\,g\,h$, d'où $\rho=\dfrac{k}{g}=\dfrac{\num{1,01e4}}{9{,}81}\approx\res{\qty{1,03e3}{kg.m^{-3}}}$. \pts{2 pts}
\item $\Delta P=\rho\,g\,h=\num{1,03e3}\times9{,}81\times20=\res{\qty{2,02e5}{Pa}}$, soit \qty{2,02}{bar} ; les relevés donnent $3{,}03-1{,}01=\qty{2,02}{bar}$ : \res{\text{accord}}. \pts{1 pt}
\item Loi de \res{\text{Boyle-Mariotte}} : $P\,V$ reste constant. $V=\dfrac{1{,}01\times6{,}0}{3{,}03}=\res{\qty{2,0}{L}}$ : le volume est divisé par trois. \pts{1,5 pt}
\item En remontant, la pression est divisée par trois : l'air des poumons, bloqué, prendrait un volume trois fois plus grand,
$6{,}0\times\dfrac{3{,}03}{1{,}01}=\res{\qty{18}{L}}$, que les poumons ne peuvent contenir : risque de \res{\text{surpression pulmonaire}}. Il doit expirer en remontant. \pts{1,5 pt}
\end{questions}
\schema[corrige]{ds4-pression}
\end{corrigeBac}

\begin{corrigeBac}{6}{Le curling}
\begin{questions}
\item Les points sont alignés : le mouvement est \res{\text{rectiligne}} ; les distances parcourues pendant \qty{0,50}{s} diminuent
(1,40 ; 1,30 ; 1,20 ; 1,10 ; 1,00 ; 0,90 m) : il est \res{\text{ralenti}}. \pts{1 pt}
\item $v_2=\dfrac{\mathrm{M_1M_3}}{2\,\Delta t}=\dfrac{3{,}90-1{,}40}{1{,}00}=\res{\qty{2,50}{m.s^{-1}}}$ ; $v_3=\dfrac{\mathrm{M_2M_4}}{2\,\Delta t}=\dfrac{5{,}00-2{,}70}{1{,}00}=\res{\qty{2,30}{m.s^{-1}}}$. \pts{1,5 pt}
\item Flèches de \qty{5,0}{cm} et \qty{4,6}{cm}, dans le sens du mouvement ; $\Delta\vec v$, de valeur \qty{0,20}{m.s^{-1}}, mesure \qty{0,4}{cm} et est
\res{\text{opposé au mouvement}} (schéma). \pts{1 pt}
\item $\sum\vec F\approx m\,\dfrac{\Delta\vec v}{\Delta t}$ : horizontale, opposée au mouvement, de valeur $19{,}0\times\dfrac{0{,}20}{0{,}50}=\res{\qty{7,6}{N}}$. \pts{1,5 pt}
\item Le poids (vertical, \qty{186}{N}), la réaction normale de la glace (verticale) et la force de frottement (horizontale, opposée
au mouvement). Les forces verticales se compensent : la force de frottement vaut \res{\qty{7,6}{N}}. \pts{1 pt}
\end{questions}
\schema[corrige]{ds4-curling}
\piege{Calculer une vitesse avec un seul intervalle : $v_i$ se calcule entre les points qui encadrent M$_i$, sur une durée $2\,\Delta t$.}
\end{corrigeBac}
